% ID: SEL-D0061-OB-E03
% ESTADO: cotejada. Segunda pasada visual y borrador compilado revisados.
% Fuente: Selectividad/sel_2010_matematicas_aplicadas/mat_apli/mat_apli_junio/mat_apli_junio/EXAMEN-ANDALUCÍA-5.pdf
% SHA-256: 5ea9fab18d8185843726e08994f39cf34f9b7dab32d86c310902a2f203dace2a
% Página: 2; opción B; ejercicio original 3.

De las 180 personas que asisten a un congreso médico, 100 son mujeres. Observando las especialidades de los congresistas, vemos que de las 60 personas que son pediatras 20 son mujeres. Se elige al azar una persona asistente al congreso.
\begin{enumerate}[label=\alph*)]
\item \textbf{(0.75 puntos)} \textquestiondown{}Cuál es la probabilidad de que sea mujer y pediatra?
\item \textbf{(0.75 puntos)} \textquestiondown{}Cuál es la probabilidad de que no sea hombre ni sea pediatra?
\item \textbf{(1 punto)} \textquestiondown{}Cuál es la probabilidad de que sea pediatra?
\end{enumerate}
